% Partie : théorie de la démonstration
% Chapitre : calcul des séquents
% Section : règles admissibles et dérivables

\documentclass[../../../include/open-logic-section]{subfiles}

\begin{document}

\olfileid[fr]{pt}{seq}{adm}

\olsection{Règles admissibles et dérivables}

De nombreux résultats de théorie de la démonstration portent
sur la possibilité de se passer d'une règle, ou utilisent
un tel résultat. Le théorème d'élimination des coupures
en est l'exemple le plus connu : tout séquent prouvable
dans un calcul muni de~\Cut{} peut aussi être prouvé
sans~\Cut. Cette propriété a un nom.

\begin{defn}
Une règle~$R$ est \emph{admissible} dans un calcul
des séquents si tout séquent prouvable dans ce calcul
augmenté de~$R$ est déjà prouvable sans~$R$.
\end{defn}

\begin{ex}\ollabel{ex:land-deriv}
La règle $\LeftR{\land}$ a des formes différentes
dans \Log{G1c} et \Log{G3c}. Dans \Log{G1c},
ses deux versions ont dans la prémisse soit
le premier conjoint $!A$ de la formule principale
$!A \land !B$, soit le second conjoint~$!B$ :
\[
\Axiom$ !A, \Gamma \fCenter \Delta$
\RightLabel{$\LeftR{\land}_1$}
\UnaryInf$ !A \land !B, \Gamma \fCenter \Delta$
\DisplayProof
\qquad
\Axiom$!B, \Gamma \fCenter \Delta$
\RightLabel{$\LeftR{\land}_2$}
\UnaryInf$!A \land !B, \Gamma \fCenter \Delta$
\DisplayProof
\]
Dans \Log{G3c}, il n'y a qu'une seule règle :
\[
\Axiom$ !A, !B, \Gamma \fCenter \Delta$
\RightLabel{$\LeftR{\land}_3$}
\UnaryInf$ !A \land !B, \Gamma \fCenter \Delta$
\DisplayProof
\]
La règle $\LeftR{\land}_3$ de \Log{G3c} est-elle
admissible dans \Log{G1c} ? Oui : toute inférence
par $\LeftR{\land}_3$ peut être remplacée par
une suite d'inférences de~\Log{G1c}. Outre
$\LeftR{\land}_1$ et $\LeftR{\land}_2$, on
emploie la règle structurelle~$\LeftR{\Contraction}$ :
\[
\Axiom$!A, !B, \Gamma \fCenter \Delta$
\RightLabel{$\LeftR{\land}_1$}
\UnaryInf$!A \land !B, !B, \Gamma \fCenter \Delta$
\RightLabel{$\LeftR{\land}_2$}
\UnaryInf$!A \land !B, !A \land !B, \Gamma \fCenter \Delta$
\RightLabel{\LeftR{\Contraction}}
\UnaryInf$!A \land !B, \Gamma \fCenter \Delta$
\DisplayProof
\]
\end{ex}

Remplacer chaque application d'une règle par des
inférences des autres règles est une façon de
prouver son admissibilité. Ce n'est pas la seule :
certaines règles admissibles ne peuvent pas
être remplacées de cette manière. Celles qui
le peuvent sont dites \emph{dérivables}.

\begin{defn}
Une règle~$R$,
\[
\AxiomC{$S_1$}
\AxiomC{$\dots$}
\AxiomC{$S_n$}
\RightLabel{$R$}
\TrinaryInfC{$S$}
\DisplayProof
\]
est \emph{dérivable} dans un système s'il existe
un schéma de preuve du séquent~$S$ utilisant
les axiomes et les règles du système, ainsi
que des séquents initiaux supplémentaires
de la forme $S_1$, \dots,~$S_n$.
\end{defn}

La preuve de l'\olref{ex:land-deriv} montre que
$\LeftR{\land}_3$ est dérivable dans~\Log{G1c}.
Elle donne un schéma de preuve de sa conclusion
à partir des prémisses de $\LeftR{\land}_3$,
prises comme séquents initiaux, par les seules
règles de \Log{G1c}. Aucun axiome de \Log{G1c}
n'y est employé.

\begin{prop}\ollabel{prop:lor-G3-adm}
Les règles $\LeftR{\land}$ et $\RightR{\lor}$
de~\Log{G3c} sont dérivables dans~\Log{G1c}.
\end{prop}

\begin{proof}
Le cas $\LeftR{\land}$ est établi dans
l'\olref{ex:land-deriv}. Le cas $\RightR{\lor}$
est laissé en exercice.
\end{proof}

\begin{prob}
Montrer que la règle $\RightR{\lor}$ de
\Log{G3c} est dérivable dans~\Log{G1c}.
\end{prob}

\begin{prob}
Supposons que $\liff$ soit un opérateur défini :
$!A \liff !B$ abrège $(!A \lif !B) \land (!B \lif !A)$.
Établir l'admissibilité des deux règles dans \Log{G1c} et \Log{G3c}.
Donner un schéma de dérivation de la règle droite.
\begin{enumerate}
\item \Axiom$\Gamma, !A, !B \fCenter \Delta$
\Axiom$\Gamma \fCenter \Delta, !A, !B$
\RightLabel{\LeftR{\liff}}
\BinaryInf$!A \liff !B, \Gamma \fCenter \Delta$
\DisplayProof
\item
\Axiom$!A, \Gamma \fCenter \Delta, !B$
\Axiom$!B, \Gamma \fCenter \Delta, !A$
\RightLabel{\RightR{\liff}}
\BinaryInf$\Gamma \fCenter \Delta, !A \liff !B$
\DisplayProof
\end{enumerate}
\end{prob}

\begin{ex}\ollabel{ex:weak-adm}
La règle $\RightR{\Weakening}$ de \Log{G1c},
\[
\Axiom$ \Gamma \fCenter \Delta$
\RightLabel{\RightR{\Weakening}}
\UnaryInf$ \Gamma \fCenter \Delta, !A$
\DisplayProof
\]
est admissible dans~\Log{G3c}.
Supposons donnée une preuve~$\pi$ de la prémisse
$\Gamma \Sequent \Delta$. Après avoir renommé
ses constantes propres pour qu'elles soient
absentes de~$!A$, on ajoute $!A$ au conséquent
de chaque séquent de~$\pi$. Les axiomes restent
des axiomes et les inférences restent correctes,
y compris les inférences quantifiées.
Le séquent final devient $\Gamma \Sequent \Delta, !A$.

La règle $\RightR{\Weakening}$ n'est pourtant
pas dérivable dans~\Log{G3c}. Supposons qu'il
existe un schéma uniforme de preuve de
$\Gamma \Sequent \Delta, !A$ utilisant les
axiomes et règles de~\Log{G3c}, avec éventuellement
$\Gamma \Sequent \Delta$ comme séquent initial
supplémentaire. En donnant à $\Gamma$ et~$\Delta$
la valeur vide, ce schéma fournirait une
preuve de~$\Sequent !A$.
Les séquents initiaux supplémentaires
$\Gamma \Sequent \Delta$ deviendraient alors
le séquent vide~$\quad \Sequent \quad$.
Or toutes les règles de~\Log{G3c} exigent
au moins une formule active dans leurs
prémisses. Un séquent vide ne peut donc
être prémisse d'aucune de ces inférences.
Il ne peut pas non plus être le séquent
final recherché. Le schéma ne peut donc
effectivement utiliser le séquent supplémentaire
$\Gamma \Sequent \Delta$.
Il donnerait une preuve de
$\Gamma \Sequent \Delta, !A$ par les seuls
axiomes et règles de~\Log{G3c}, pour toute
valeur de ces paramètres. C'est impossible :
par correction, \Log{G3c} ne prouve que
des séquents valides. Pour $\Gamma = \Delta = \emptyset$
et $!A \ident \lfalse$, il devrait notamment
prouver $\quad \Sequent \lfalse$.
\end{ex}

On peut donner une preuve par récurrence
de l'admissibilité de l'affaiblissement.
L'argument précédent établit même davantage :
l'ajout de $!A$ à tous les conséquents
de~$\pi$ conserve la structure de l'arbre,
et donc sa hauteur. Cette propriété mérite
une définition.

\begin{defn}
Une règle~$R$,
\[
\AxiomC{$S_1$}
\AxiomC{$\dots$}
\AxiomC{$S_n$}
\RightLabel{$R$}
\TrinaryInfC{$S$}
\DisplayProof
\]
est \emph{admissible sans augmentation de hauteur},
ou \emph{hp-admissible}, si $\Proves[h_i] S_i$
pour $i = 1$, \dots,~$n$ entraîne $\Proves[m] S$,
où $m = \max(h_1, \dots, h_n)$.
\end{defn}

Il suffit donc que, pour des preuves~$\pi_i$
des prémisses~$S_i$ de hauteurs
$h_i = \pheight{\pi_i}$, on puisse construire
une preuve~$\pi$ de la conclusion dont la
hauteur~$\pheight{\pi}$ ne dépasse pas
le maximum des~$h_i$.
Pour une seule prémisse~$S_1$, la condition
est $\pheight{\pi} \le \pheight{\pi_1}$.
Pour $\RightR{\Weakening}$ et $\LeftR{\Weakening}$,
la construction donne même
$\pheight{\pi} = \pheight{\pi_1}$,
mais cette égalité n'est pas exigée
par la définition générale.

\begin{prop}\ollabel{prop:weak-G3c-adm}
Les règles $\RightR{\Weakening}$ et $\LeftR{\Weakening}$
de \Log{G1c} sont admissibles sans augmentation
de hauteur dans~\Log{G3c}.
\end{prop}

\begin{proof}
Supposons $\Log{G3c} \Proves \Gamma \Sequent \Delta$
avec une preuve~$\pi$ de hauteur~$\pheight{\pi} = n$.
Après renommage de ses constantes propres,
qui conserve sa hauteur et son séquent final,
on les suppose absentes de la formule ajoutée.
Il faut construire une preuve~$\pi'$ de
$\Gamma \Sequent \Delta, !A$ dans \Log{G3c},
avec $\pheight{\pi'} = n$.
On raisonne par récurrence sur~$n$.

Si $n=0$, $\pi$ est l'axiome
$\Gamma \Sequent \Delta$ : soit une même
formule atomique~$!C$ figure dans $\Gamma$
et~$\Delta$, soit l'antécédent contient
l'absurde. Dans les deux cas,
$\Gamma \Sequent \Delta, !A$ reste un axiome.

Si $n = \pheight{\pi} >0$, on distingue
les dernières inférences. Voici le cas
$\RightR{\land}$. On a $\Delta = \Delta', !B \land !C$,
et les sous-preuves immédiates $\pi_1$ et~$\pi_2$
se terminent respectivement par
$\Gamma \Sequent \Delta', !B$ et
$\Gamma \Sequent \Delta', !C$.
La preuve $\pi$ a donc la forme suivante :
\[
  \AxiomC{}
  \RightLabel{$\pi_1$}
  \Deduce$\Gamma \fCenter \Delta', !B$
  \AxiomC{}
  \RightLabel{$\pi_2$}
  \Deduce$\Gamma \fCenter \Delta', !C$
  \RightLabel{\RightR{\land}}
  \BinaryInf$\Gamma \fCenter \Delta', !B \land !C$
  \DisplayProof
\]
Comme $n = \max(\pheight{\pi_1}, \pheight{\pi_2}) + 1$,
on a $\pheight{\pi_1} < n$ et $\pheight{\pi_2} < n$.
L'hypothèse de récurrence fournit des preuves
$\pi_1'$ et~$\pi_2'$ de
$\Gamma \Sequent \Delta', !B, !A$ et
$\Gamma \Sequent \Delta', !C, !A$.
Par $\RightR{\land}$, on obtient une preuve~$\pi'$
de $\Gamma \Sequent \Delta', !B \land !C, !A$ :
\[
  \AxiomC{}
  \RightLabel{$\pi_1'$}
  \Deduce$\Gamma \fCenter \Delta', !B, !A$
  \AxiomC{}
  \RightLabel{$\pi_2'$}
  \Deduce$\Gamma \fCenter \Delta', !C, !A$
  \RightLabel{\RightR{\land}}
  \BinaryInf$\Gamma \fCenter \Delta', !B \land !C, !A$
  \DisplayProof
\]
Sa hauteur est
$\pheight{\pi'} = \max(\pheight{\pi_1'}, \pheight{\pi_2'}) + 1
= \max(\pheight{\pi_1}, \pheight{\pi_2}) + 1 = \pheight{\pi}$.
Les autres règles se traitent de la même manière :
on garde l'inférence et on ajoute la formule
au contexte des prémisses et de la conclusion.
Pour les règles à constante propre, le renommage
préalable assure toujours leur condition.
L'affaiblissement à gauche se prouve de même.
\end{proof}

La \olref{prop:weak-G3c-adm}
sera souvent utilisée. On note
$\pi[\Pi \Sequent \Lambda]$ le résultat
de l'application répétée de l'affaiblissement
admissible à une preuve~$\pi$ de
$\Gamma \Sequent \Delta$ : on ajoute
à gauche les formules de~$\Pi$, et à droite
celles de~$\Lambda$. On obtient ainsi une
preuve de $\Gamma, \Pi \Sequent \Delta, \Lambda$.

Le renommage des constantes propres s'applique aussi
lorsque les formules ajoutées sont closes. Une telle
formule peut contenir une constante employée comme
%
%
%
%
%
%
%
%
%
%
%
constante propre dans une inférence intérieure.

\end{document}
